% Source: Solutions/Karan-Gangadharan-PS7-Q1-4.tex (LaTeX source). "Question 2(a)" renumbered to "Problem 7.2(a)".
\submittedby{Karan Gangadharan}

\begin{proof}
Let
\[ \pi_1\colon\R\times\R\to\R \]
be the projection onto the first coordinate, defined by $\pi_1(x,y)=x$. Let
\[ A=\{(x,y)\in\R\times\R : x\ge0\text{ or }y=0\}. \]
Define $g\colon A\to\R$ by restricting $\pi_1$ to $A$. Thus, $g(x,y)=x$.

We shall show that $g$ is a quotient map but is neither open nor closed.

\medskip
\textbf{(i) $g$ is a quotient map:}
Since $g$ is the restriction of the continuous projection map $\pi_1$, it is continuous. Also,
$g$ is surjective, because for every $x\in\R$,
\[ (x,0)\in A\quad\text{and}\quad g(x,0)=x. \]
Now, consider the map $f\colon\R\to A$ defined by $f(x)=(x,0)$. Clearly, $f$ is continuous and
\[ (g\circ f)(x)=g(x,0)=x. \]
Therefore, $g\circ f=I_\R$. By Problem 7.2(a), since $g$ and $f$ are continuous and
$g\circ f=I_\R$, it follows that
\[ \boxed{g\text{ is a quotient map.}} \]

\medskip
\textbf{(ii) $g$ is not an open map:}
Consider the set
\[ U=A\cap\bigl((-1,1)\times(1/2,3/2)\bigr). \]
Since $(-1,1)\times(1/2,3/2)$ is open in $\R^2$, $U$ is open in $A$. Now, by the definition of
$A$, if $(x,y)\in U$, then $x\ge0$, since $y\neq0$. Consequently,
\[ g(U)=[0,1). \]
But $[0,1)$ is not open in $\R$. Thus, $g$ does not map every open subset of $A$ onto an open
subset of $\R$. Hence,
\[ \boxed{g\text{ is not an open map.}} \]

\medskip
\textbf{(iii) $g$ is not a closed map:}
Consider the subset
\[ F=\left\{\left(\frac1n,n\right) : n\in\N\right\}. \]
Clearly, $F\subseteq A$, since $1/n>0$ for every $n\in\N$. We claim that $F$ is closed in $A$.
Indeed, $F$ is closed in $\R^2$, because the second coordinates of its points tend to infinity,
so it has no accumulation point in $\R^2$. Therefore, $F$ is also closed in the subspace $A$.
Now,
\[ g(F)=g\left(\left\{\left(\frac1n,n\right) : n\in\N\right\}\right)
       =\left\{\frac1n : n\in\N\right\}. \]
But
\[ \frac1n\longrightarrow0\quad\text{as }n\longrightarrow\infty, \]
and $0\notin g(F)$. Hence, $g(F)$ is not closed in $\R$. Therefore, $g$ does not map every
closed subset of $A$ onto a closed subset of $\R$. Thus,
\[ \boxed{g\text{ is not a closed map.}} \]

Combining (i), (ii) and (iii), we conclude that
\[ \boxed{g\colon A\to\R\text{ is a quotient map which is neither open nor closed.}} \qedhere \]
\end{proof}
